% Part: second-order-logic
% Chapter: syntax-and-semantics
% Section: size-of-domain

\documentclass[../../../include/open-logic-section]{subfiles}

\begin{document}

\olfileid{sol}{syn}{siz}

\olsection{અનંત અને \usetoken{S}{enumerable}
  \usetoken{P}{domain}નું વર્ણન}

ગણ~$M$ (ડેડેકિન્ડના અર્થમાં) અનંત ત્યારે અને માત્ર ત્યારે જ હોય જ્યારે
એવું !!a{injective} વિધેય~$f\colon M
\to M$ હોય જે !!{surjective} ન હોય, એટલે $\ran{f} \neq M$. પ્રથમ-ક્રમ
તર્કશાસ્ત્રમાં એકસ્થાની !!{function}~$f$ લઈ શકાય અને
!!a{structure}~$\Struct{M}$માં તેને અપાયેલું વિધેય~$\Assign{f}{M}$
!!{injective} છે તથા $\ran{f} \neq
\Domain{M}$ છે એમ કહી શકાય:
\[
\lforall[x][\lforall[y][(\eq[f(x)][f(y)] \to \eq[x][y])]] \land
\lexists[y][\lforall[x][\eq/[y][f(x)]]].
\]
જો~$\Struct{M}$ આ !!{sentence}ને સંતોષે, તો
$\Assign{f}{M}:
\Domain{M} \to \Domain{M}$ !!{injective} છે અને તેથી~$\Domain{M}$
અનંત હોવું જોઈએ. જો~$\Domain{M}$ અનંત હોય અને તેથી આવું વિધેય
અસ્તિત્વમાં હોય, તો~$\Assign{f}{M}$ને તે વિધેય લઈએ; ત્યારે
$\Struct{M}$ આ !!{sentence}ને સંતોષશે. પરંતુ આ માટે આપણા હેતુનો
અતાર્કિક સંકેત~$f$ ભાષામાં હોવો જરૂરી છે. દ્વિતીય-ક્રમ તર્કશાસ્ત્રમાં
આવું વિધેય \emph{અસ્તિત્વમાં છે} એમ સીધું કહી શકાય છે. ત્યારે~$f$ની
જરૂર રહેતી નથી અને શુદ્ધ દ્વિતીય-ક્રમ તર્કશાસ્ત્રમાં આ !!{sentence}
મળે છે:
\[
\fn{Inf} \ident \lexists[u][(\lforall[x][\lforall[y][(\eq[u(x)][u(y)]
      \lif \eq[x][y])]] \land \lexists[y][\lforall[x][\eq/[y][u(x)]]])].
\]
$\Sat{M}{\fn{Inf}}$ ત્યારે અને માત્ર ત્યારે જ જ્યારે~$\Domain{M}$
અનંત હોય. પછી $\fn{Fin} \ident \lnot \fn{Inf}$ વ્યાખ્યાયિત કરી
શકાય; $\Sat{M}{\fn{Fin}}$ ત્યારે અને માત્ર ત્યારે જ જ્યારે
$\Domain{M}$ સાન્ત હોય. શુદ્ધ પ્રથમ-ક્રમ તર્કશાસ્ત્રનું એક પણ એકલું
!!{sentence} !!{domain} અનંત છે એવું વ્યક્ત કરી શકતું નથી, જોકે
આવાં વાક્યોનો અનંત ગણ તે વ્યક્ત કરી શકે છે. શુદ્ધ પ્રથમ-ક્રમ
તર્કશાસ્ત્રનાં !!{sentence}sનો એવો કોઈ ગણ નથી જે કોઈ
!!a{structure}માં ત્યારે અને માત્ર ત્યારે જ સંતોષાય જ્યારે તેનું
વ્યાપકક્ષેત્ર સાન્ત હોય.

\begin{prop}
$\Sat{M}{\fn{Inf}}$ ત્યારે અને માત્ર ત્યારે જ જ્યારે~$\Domain{M}$
અનંત હોય.
\end{prop}

\begin{proof}
$\Sat{M}{\fn{Inf}}$ ત્યારે અને માત્ર ત્યારે જ જ્યારે અમુક~$s$
માટે
  $\Sat{M}{\lforall[x][\lforall[y][(\eq[u(x)][u(y)] \lif \eq[x][y])]]
    \land \lexists[y][\lforall[x][\eq/[y][u(x)]]]}[s]$.
  એવું હોય તો~$s(u)$ !!a{injective} વિધેય છે અને અમુક $y
  \in \Domain{M}$ એ~$s(u)$ના મૂલ્યસમૂહમાં નથી. ઊલટું, જો એવું
  !!a{injective} $f\colon \Domain{M} \to \Domain{M}$ હોય કે
  $\ran{f} \neq \Domain{M}$, તો~$s(u) = f$ એવી ચલ-નિયુક્તિ જરૂરી
  શરત પૂરી કરે છે.
\end{proof}

અહીં અરિક્ત ગણ~$M$ને !!{enumerable} કહીએ છીએ જો તેના !!{element}sની
આવી પરિગણના હોય:
\[
m_0, m_1, m_2, \dots
\]
તેમાં પુનરાવર્તન નથી, પરંતુ પરિગણના સાન્ત હોઈ શકે છે.
\sourcecorrection{OLSOL-004}{સ્થિર અંગ્રેજી મૂળ અહીં મનસ્વી ગણને
પરિગણનીય કહે છે, પણ તરત પછી $m_0$ અને $z\in M$ માગે છે.
રિક્ત ગણ માટે આ શક્ય નથી. સંરચનાનાં વ્યાપકક્ષેત્રો અરિક્ત હોય છે,
અને અહીં જરૂરી અરિક્તતાની શરત સ્પષ્ટ કરી છે.}
આવી પરિગણના ત્યારે અને માત્ર ત્યારે જ મળે જ્યારે !!a{element}
$z \in M$ અને વિધેય $f\colon M \to M$ એવાં હોય કે $z$, $f(z)$,
$f(f(z))$, \dots, ના જુદા જુદા ઘટકો જ~$M$ના બધા !!{element}s હોય.
કારણ કે જો પરિગણના હોય, તો $z = m_0$ અને $f(m_k) = m_{k+1}$
(અથવા~$m_k$ પરિગણનાનો છેલ્લો !!{element} હોય તો
$f(m_k) = m_k$) જરૂરી !!{element} અને વિધેય આપે છે. બીજી બાજુ,
આવાં~$z$ અને~$f$ હોય તો $z$, $f(z)$, $f(f(z))$, \dots, માંથી
પુનરાવર્તિત ઘટકો કાઢતાં~$M$ની પરિગણના મળે છે અને~$M$
!!{enumerable} છે. આ~$z$ અને~$f$નું અસ્તિત્વ દ્વિતીય-ક્રમ
તર્કશાસ્ત્રમાં વ્યક્ત કરીને એવું !!{sentence} બનાવી શકાય છે જે
!!a{structure}માં ત્યારે અને માત્ર ત્યારે જ સાચું હોય જ્યારે તે
!!{structure} !!{enumerable} હોય:
\[
\fn{Count} \ident
\lexists[z][\lexists[u][\lforall[X][((X(z) \land
      \lforall[x][(X(x) \lif X(u(x)))]) \lif \lforall[x][X(x)])]]]
\]
\sourcecorrection{OLSOL-005}{સ્થિર અંગ્રેજી મૂળમાં પરિગણના
પુનરાવર્તન વિનાની કહેવાઈ છે, પણ સાન્ત પરિગણનાના છેલ્લા ઘટક પર
વિધેય સ્થિર થાય ત્યારે તેની આગળની કક્ષા એ ઘટકને વારંવાર આપે છે.
અહીં કક્ષાના જુદા ઘટકો પરિગણના બનાવે છે એમ સ્પષ્ટ કર્યું છે;
નીચેની સાબિતીમાં પણ છેલ્લા ઘટકનું વિધેય-મૂલ્ય સ્પષ્ટ કરાયું છે.}

\begin{prop}
$\Sat{M}{\fn{Count}}$ ત્યારે અને માત્ર ત્યારે જ જ્યારે~$\Domain{M}$
!!{enumerable} હોય.
\end{prop}

\begin{proof}
માનો કે~$\Domain{M}$ !!{enumerable} છે અને $m_0$, $m_1$, \dots,
તેની પરિગણના છે. પુનરાવર્તનો દૂર કરીને કોઈ~$m_k$ બે વાર ન આવે તેની
ખાતરી કરી શકાય છે. $f(m_k) = m_{k+1}$ વ્યાખ્યાયિત કરો (સાન્ત
પરિગણનાના છેલ્લા ઘટકને પોતાને જ મોકલો) અને $s(z) = m_0$ તથા
$s(u) = f$ લો. આપણે બતાવીએ કે
\[
\Sat{M}{\lforall[X][((X(z) \land \lforall[x][(X(x) \lif X(u(x)))])
    \lif \lforall[x][X(x)])]}[s]
\]
માનો કે $M \subseteq \Domain{M}$ મનસ્વી છે. વધુમાં માનો કે
$\Sat{M}{(X(z) \land \lforall[x][(X(x) \lif
X(u(x)))])}[\Subst{s}{M}{X}]$. ત્યારે
$\Subst{s}{M}{X}(z) \in M$ અને જ્યારે પણ $x \in M$, ત્યારે
$(\Subst{s}{M}{X}(u))(x) \in M$ પણ. બીજા શબ્દોમાં,
$\varAssign{\Subst{s}{M}{X}}{s}{X}$ હોવાથી, $m_0 \in M$ અને
$x \in M$ હોય તો $f(x) \in M$. તેથી $m_0 \in M$,
$m_1 = f(m_0) \in M$, $m_2 = f(f(m_0)) \in M$ વગેરે. આમ,
$M = \Domain{M}$ અને તેથી
$\Sat{M}{\lforall[x][X(x)]}[\Subst{s}{M}{X}]$. $M \subseteq
\Domain{M}$ મનસ્વી હોવાથી સાબિતી પૂરી થાય છે:
$\Sat{M}{\fn{Count}}$.

હવે માનો કે $\Sat{M}{\fn{Count}}$, એટલે અમુક~$s$ માટે
\[
\Sat{M}{\lforall[X][((X(z) \land \lforall[x][(X(x) \lif X(u(x)))])
    \lif \lforall[x][X(x)])]}[s]
\]
$m = s(z)$ અને $f = s(u)$ લો તથા
$M = \{m,
f(m), f(f(m)), \dots\}$ વિચારો. પુનરાવર્તિત ઘટકો કાઢીએ તો
આ~$M$ !!{enumerable} છે. ત્યારે ધારણા પરથી
\[
\Sat{M}{(X(z) \land \lforall[x][(X(x) \lif X(u(x)))] )
    \lif \lforall[x][X(x)]}[\Subst{s}{M}{X}]
\]
મળે છે. વળી $M \ni m =
\Subst{s}{M}{X}(z)$ હોવાથી $\Sat{M}{X(z)}[\Subst{s}{M}{X}]$,
અને $x \in M$ હોય ત્યારે $f(x) \in
M$ હોવાથી $\Sat{M}{\lforall[x][(X(x) \lif
X(u(x)))]}[\Subst{s}{M}{X}]$ પણ મળે છે. આમ, પૂર્વપક્ષ અને
શરતી વિધાન બંને સંતોષાય છે, તેથી ઉત્તરપક્ષ પણ સંતોષાવો જોઈએ:
$\Sat{M}{\lforall[x][X(x)]}[\Subst{s}{M}{X}]$. પરંતુ તેનો અર્થ
$M = \Domain{M}$; અને~$M$ વ્યાખ્યા મુજબ પરિગણનીય હોવાથી
$\Domain{M}$ પણ !!{enumerable} છે.
\end{proof}

\begin{prob}
!!{sentence}~$\fn{Inf} \land \fn{Count}$ બધા અને માત્ર
!!{denumerable} વ્યાપકક્ષેત્રોમાં સાચું છે. $\fn{Count}$ની વ્યાખ્યા
એવી બદલો કે નવું !!{sentence} વ્યાપકક્ષેત્ર !!{denumerable} છે એવું
સીધું વ્યક્ત કરે, અને તે કરે છે એમ સાબિત કરો.
\end{prob}

\end{document}
